% !TeX root = ../main.tex

\chapter{Composition}

\section{Composition by Powers}
% =================================================

Consider

\[
\phi_n(x)=b^{x^n}.
\]

Replace \(x\) by \(x^m\). Then

\[
\phi_n(x^m)
=
b^{(x^m)^n}
=
b^{x^{mn}}.
\]

Therefore

\[
\boxed{
\phi_n(x^m)=\phi_{mn}(x).
}
\]

For example,

\[
\phi_2(x^3)
=
b^{(x^3)^2}
=
b^{x^6}
=
\phi_6(x),
\]

while

\[
\phi_3(x^2)
=
b^{(x^2)^3}
=
b^{x^6}
=
\phi_6(x).
\]

Hence

\[
\boxed{
\phi_2(x^3)
=
\phi_3(x^2)
=
\phi_6(x).
}
\]

More generally,

\[
\boxed{
\phi_n(x^m)
=
\phi_m(x^n)
=
\phi_{mn}(x).
}
\]

Thus multiplication of natural-number indices corresponds to
composition with power maps.

% =================================================
\section{Composition Operators}
% =================================================

The previous phenomenon becomes clearer if we introduce an operator.

For every \(m\in\N\), define

\[
\boxed{
T_mF(x)=F(x^m).
}
\]

Then, when \(T_m\) acts on one of our basic blocks,

\[
T_m\phi_n(x)
=
\phi_n(x^m).
\]

Therefore

\[
\boxed{
T_m\phi_n=\phi_{mn}.
}
\]

Now consider two such operators:

\[
T_mT_nF(x).
\]

We first apply \(T_n\):

\[
T_nF(x)=F(x^n).
\]

Then

\[
T_mT_nF(x)
=
F((x^m)^n)
=
F(x^{mn}).
\]

Hence

\[
\boxed{
T_mT_n=T_{mn}.
}
\]

Since multiplication of natural numbers is commutative,

\[
mn=nm,
\]

we also have

\[
\boxed{
T_mT_n=T_nT_m.
}
\]

Finally,

\[
T_1F(x)=F(x),
\]

so

\[
\boxed{
T_1=I.
}
\]

We therefore obtain the same algebraic laws as the multiplicative
monoid

\[
(\N,\cdot).
\]



\noindent\textbf{Concrete examples.}
\begin{itemize}
  \item If \(F(x)=\phi_2(x)+3\phi_5(x)=b^{x^2}+3b^{x^5}\), then
  

\[
  T_3F(x)=F(x^3)=b^{(x^3)^2}+3b^{(x^3)^5}=b^{x^6}+3b^{x^{15}}=\phi_6(x)+3\phi_{15}(x).
  \]


  \item If \(G(x)=\sum_{n\ge1} c_n\phi_n(x)\), then
  

\[
  T_mG(x)=\sum_{n\ge1} c_n\phi_n(x^m)=\sum_{n\ge1} c_n\phi_{mn}(x),
  \]


  so the coefficient at index \(n\) is transported to index \(mn\).
\end{itemize}

\medskip

\noindent These formulas make explicit the correspondence between functional composition with power maps and multiplication of natural-number indices.

% =================================================
